% Source: https://artofproblemsolving.com/community/c3h1443160_committee_problem_from_a_generalization
% Stable identity: aops-topic-1443160-q01
\textbf{(a)} The committee is any $n$-element subset of the $2n$ members, so there are $\dbinom{2n}{n}$ committees.

\textbf{(b)} A committee with exactly $k$ boys has $n-k$ girls. The boys can be chosen in $\dbinom nk$ ways and the girls in $\dbinom{n}{n-k}=\dbinom nk$ ways, so there are $\dbinom nk^2$ such committees. Every committee has some number $k$ of boys with $0\le k\le n$, and different values of $k$ give different committees, so
\[
\sum_{k=0}^{n}\dbinom nk^2=\dbinom{2n}{n}. \qquad\blacksquare
\]
For example, with $n=3$: $1+9+9+1=20=\dbinom63$.
